% ===========================================================================
% CTF 网络安全入门：Web 专攻版
% 重构版讲义 —— 面向零基础初学者，知识点铺垫 + 标准题面卡 + 逐题图解带做
% 所有截图均为真实执行记录（终端实录截图 / 本地靶场真实网页截图 / 原理示意图）
% ===========================================================================
\documentclass[UTF8,12pt,oneside]{ctexbook}
% handout-common.tex —— 五本讲义共享的导言区与排版宏
% 所有 handout-*.tex 通过 % handout-common.tex —— 五本讲义共享的导言区与排版宏
% 所有 handout-*.tex 通过 % handout-common.tex —— 五本讲义共享的导言区与排版宏
% 所有 handout-*.tex 通过 \input{handout-common} 引入本文件。
% 【重要】此文件为统一规范，任何单本讲义不得修改它。
% ===========================================================================
\usepackage[a4paper,left=18mm,right=18mm,top=18mm,bottom=18mm,
  includeheadfoot,headheight=15pt,headsep=4mm,footskip=10mm]{geometry}
\usepackage{amsmath,amssymb,booktabs,longtable,array,enumitem,listings,xcolor,hyperref,fancyhdr,graphicx}
\graphicspath{{figures/}}
\hypersetup{hidelinks,pdfauthor={CTF 培训资料}}

\setlength{\parindent}{2em}
\setlength{\parskip}{0.25em}
\emergencystretch=3em
\setlist{itemsep=0.15em,topsep=0.25em}
\lstset{basicstyle=\ttfamily\small,breaklines=true,columns=fullflexible,keepspaces=true,
  showstringspaces=false,frame=single,framerule=0.3pt,xleftmargin=0.4em,xrightmargin=0.4em}

% ---- 页眉页脚 -------------------------------------------------------------
\pagestyle{fancy}
\fancyhf{}
\fancyhead[L]{CTF 网络安全入门}
\fancyhead[R]{\thehandoutmark}
\fancyfoot[C]{\thepage}
\renewcommand{\headrulewidth}{0.3pt}
\newcommand{\thehandoutmark}{讲义}
% 用法（写在导言区）：\handoutmark{Web 专攻版}
\newcommand{\handoutmark}[1]{\renewcommand{\thehandoutmark}{#1}}
\newcommand{\booktitle}[1]{\hypersetup{pdftitle={#1}}}

% ---- 正文辅助宏 -----------------------------------------------------------
\newcommand{\goal}[1]{\noindent\textbf{本章目标：}#1\par}
\newcommand{\code}[1]{\texttt{#1}}
\newcommand{\kw}[1]{\textbf{#1}}
\newcommand{\flagbox}[1]{\par\noindent{\setlength{\fboxsep}{4pt}\fbox{\textbf{flag：}\texttt{#1}}}\par}
% 题面卡中的字段行：\field{附件}{...}
\newcommand{\field}[2]{\par\noindent\textbf{#1：}#2\par}

% ---- 信息框（一句话记忆 / 常见误区 / 排错指南 / 小技巧）--------------------
\newsavebox{\ibox}
\newenvironment{infobox}[1]
  {\par\medskip\begin{lrbox}{\ibox}\begin{minipage}{0.95\linewidth}\textbf{#1}\par\smallskip\hrule\smallskip}
  {\end{minipage}\end{lrbox}\noindent\fbox{\usebox{\ibox}}\par\medskip}
\newcommand{\memory}[1]{\begin{infobox}{一句话记忆}#1\end{infobox}}
\newcommand{\pitfall}[1]{\begin{infobox}{常见误区}#1\end{infobox}}
\newcommand{\debugtip}[1]{\begin{infobox}{排错指南}#1\end{infobox}}
\newcommand{\tip}[1]{\begin{infobox}{小技巧}#1\end{infobox}}

% ---- 题面卡 ---------------------------------------------------------------
% \begin{challenge}{题 R1：一句话标题} ... \field{...}{...} ... \end{challenge}
\newenvironment{challenge}[1]
  {\par\medskip\begin{lrbox}{\ibox}\begin{minipage}{0.95\linewidth}{\large\bfseries #1}\par\smallskip\hrule\smallskip}
  {\end{minipage}\end{lrbox}\noindent\fbox{\usebox{\ibox}}\par\medskip}

% ---- 插图 -----------------------------------------------------------------
% \shot[宽度比例]{文件名}{图注}{标签}   引用：\figref{标签}
\newcommand{\shot}[4][0.96]{%
  \begin{figure}[htbp]\centering
  \includegraphics[width=#1\linewidth]{#2}%
  \caption{#3}\label{fig:#4}\end{figure}}
\newcommand{\figref}[1]{图~\ref{fig:#1}}

\endinput
 引入本文件。
% 【重要】此文件为统一规范，任何单本讲义不得修改它。
% ===========================================================================
\usepackage[a4paper,left=18mm,right=18mm,top=18mm,bottom=18mm,
  includeheadfoot,headheight=15pt,headsep=4mm,footskip=10mm]{geometry}
\usepackage{amsmath,amssymb,booktabs,longtable,array,enumitem,listings,xcolor,hyperref,fancyhdr,graphicx}
\graphicspath{{figures/}}
\hypersetup{hidelinks,pdfauthor={CTF 培训资料}}

\setlength{\parindent}{2em}
\setlength{\parskip}{0.25em}
\emergencystretch=3em
\setlist{itemsep=0.15em,topsep=0.25em}
\lstset{basicstyle=\ttfamily\small,breaklines=true,columns=fullflexible,keepspaces=true,
  showstringspaces=false,frame=single,framerule=0.3pt,xleftmargin=0.4em,xrightmargin=0.4em}

% ---- 页眉页脚 -------------------------------------------------------------
\pagestyle{fancy}
\fancyhf{}
\fancyhead[L]{CTF 网络安全入门}
\fancyhead[R]{\thehandoutmark}
\fancyfoot[C]{\thepage}
\renewcommand{\headrulewidth}{0.3pt}
\newcommand{\thehandoutmark}{讲义}
% 用法（写在导言区）：\handoutmark{Web 专攻版}
\newcommand{\handoutmark}[1]{\renewcommand{\thehandoutmark}{#1}}
\newcommand{\booktitle}[1]{\hypersetup{pdftitle={#1}}}

% ---- 正文辅助宏 -----------------------------------------------------------
\newcommand{\goal}[1]{\noindent\textbf{本章目标：}#1\par}
\newcommand{\code}[1]{\texttt{#1}}
\newcommand{\kw}[1]{\textbf{#1}}
\newcommand{\flagbox}[1]{\par\noindent{\setlength{\fboxsep}{4pt}\fbox{\textbf{flag：}\texttt{#1}}}\par}
% 题面卡中的字段行：\field{附件}{...}
\newcommand{\field}[2]{\par\noindent\textbf{#1：}#2\par}

% ---- 信息框（一句话记忆 / 常见误区 / 排错指南 / 小技巧）--------------------
\newsavebox{\ibox}
\newenvironment{infobox}[1]
  {\par\medskip\begin{lrbox}{\ibox}\begin{minipage}{0.95\linewidth}\textbf{#1}\par\smallskip\hrule\smallskip}
  {\end{minipage}\end{lrbox}\noindent\fbox{\usebox{\ibox}}\par\medskip}
\newcommand{\memory}[1]{\begin{infobox}{一句话记忆}#1\end{infobox}}
\newcommand{\pitfall}[1]{\begin{infobox}{常见误区}#1\end{infobox}}
\newcommand{\debugtip}[1]{\begin{infobox}{排错指南}#1\end{infobox}}
\newcommand{\tip}[1]{\begin{infobox}{小技巧}#1\end{infobox}}

% ---- 题面卡 ---------------------------------------------------------------
% \begin{challenge}{题 R1：一句话标题} ... \field{...}{...} ... \end{challenge}
\newenvironment{challenge}[1]
  {\par\medskip\begin{lrbox}{\ibox}\begin{minipage}{0.95\linewidth}{\large\bfseries #1}\par\smallskip\hrule\smallskip}
  {\end{minipage}\end{lrbox}\noindent\fbox{\usebox{\ibox}}\par\medskip}

% ---- 插图 -----------------------------------------------------------------
% \shot[宽度比例]{文件名}{图注}{标签}   引用：\figref{标签}
\newcommand{\shot}[4][0.96]{%
  \begin{figure}[htbp]\centering
  \includegraphics[width=#1\linewidth]{#2}%
  \caption{#3}\label{fig:#4}\end{figure}}
\newcommand{\figref}[1]{图~\ref{fig:#1}}

\endinput
 引入本文件。
% 【重要】此文件为统一规范，任何单本讲义不得修改它。
% ===========================================================================
\usepackage[a4paper,left=18mm,right=18mm,top=18mm,bottom=18mm,
  includeheadfoot,headheight=15pt,headsep=4mm,footskip=10mm]{geometry}
\usepackage{amsmath,amssymb,booktabs,longtable,array,enumitem,listings,xcolor,hyperref,fancyhdr,graphicx}
\graphicspath{{figures/}}
\hypersetup{hidelinks,pdfauthor={CTF 培训资料}}

\setlength{\parindent}{2em}
\setlength{\parskip}{0.25em}
\emergencystretch=3em
\setlist{itemsep=0.15em,topsep=0.25em}
\lstset{basicstyle=\ttfamily\small,breaklines=true,columns=fullflexible,keepspaces=true,
  showstringspaces=false,frame=single,framerule=0.3pt,xleftmargin=0.4em,xrightmargin=0.4em}

% ---- 页眉页脚 -------------------------------------------------------------
\pagestyle{fancy}
\fancyhf{}
\fancyhead[L]{CTF 网络安全入门}
\fancyhead[R]{\thehandoutmark}
\fancyfoot[C]{\thepage}
\renewcommand{\headrulewidth}{0.3pt}
\newcommand{\thehandoutmark}{讲义}
% 用法（写在导言区）：\handoutmark{Web 专攻版}
\newcommand{\handoutmark}[1]{\renewcommand{\thehandoutmark}{#1}}
\newcommand{\booktitle}[1]{\hypersetup{pdftitle={#1}}}

% ---- 正文辅助宏 -----------------------------------------------------------
\newcommand{\goal}[1]{\noindent\textbf{本章目标：}#1\par}
\newcommand{\code}[1]{\texttt{#1}}
\newcommand{\kw}[1]{\textbf{#1}}
\newcommand{\flagbox}[1]{\par\noindent{\setlength{\fboxsep}{4pt}\fbox{\textbf{flag：}\texttt{#1}}}\par}
% 题面卡中的字段行：\field{附件}{...}
\newcommand{\field}[2]{\par\noindent\textbf{#1：}#2\par}

% ---- 信息框（一句话记忆 / 常见误区 / 排错指南 / 小技巧）--------------------
\newsavebox{\ibox}
\newenvironment{infobox}[1]
  {\par\medskip\begin{lrbox}{\ibox}\begin{minipage}{0.95\linewidth}\textbf{#1}\par\smallskip\hrule\smallskip}
  {\end{minipage}\end{lrbox}\noindent\fbox{\usebox{\ibox}}\par\medskip}
\newcommand{\memory}[1]{\begin{infobox}{一句话记忆}#1\end{infobox}}
\newcommand{\pitfall}[1]{\begin{infobox}{常见误区}#1\end{infobox}}
\newcommand{\debugtip}[1]{\begin{infobox}{排错指南}#1\end{infobox}}
\newcommand{\tip}[1]{\begin{infobox}{小技巧}#1\end{infobox}}

% ---- 题面卡 ---------------------------------------------------------------
% \begin{challenge}{题 R1：一句话标题} ... \field{...}{...} ... \end{challenge}
\newenvironment{challenge}[1]
  {\par\medskip\begin{lrbox}{\ibox}\begin{minipage}{0.95\linewidth}{\large\bfseries #1}\par\smallskip\hrule\smallskip}
  {\end{minipage}\end{lrbox}\noindent\fbox{\usebox{\ibox}}\par\medskip}

% ---- 插图 -----------------------------------------------------------------
% \shot[宽度比例]{文件名}{图注}{标签}   引用：\figref{标签}
\newcommand{\shot}[4][0.96]{%
  \begin{figure}[htbp]\centering
  \includegraphics[width=#1\linewidth]{#2}%
  \caption{#3}\label{fig:#4}\end{figure}}
\newcommand{\figref}[1]{图~\ref{fig:#1}}

\endinput

% Keep instructional figures beside the step that introduces them.
\usepackage{float}
\renewcommand{\shot}[4][0.92]{%
  \begin{figure}[H]\centering
  \includegraphics[width=#1\linewidth]{#2}%
  \caption{#3}\label{fig:#4}\end{figure}}
% 字形修正：①②③④ 等序号与 ★☆、↔、希腊字母在西文字体中缺字形，声明为 CJK 字符类，
% 改由中文字体渲染，避免这些字符在 PDF 里被吞掉。
\xeCJKDeclareCharClass{CJK}{"2460 -> "24FF, "2194 -> "2195, "2605 -> "2606, "03B1 -> "03C9}
\handoutmark{Web 专攻版}
\booktitle{CTF 网络安全入门：Web 专攻版}

\begin{document}
\frontmatter
\begin{titlepage}
\centering
\vspace*{28mm}
{\Huge\bfseries CTF 网络安全入门\par}
\vspace{9mm}
{\LARGE Web 专攻版\par}
\vspace{5mm}
{\large 题目驱动 · 请求、权限与注入\par}
\vspace{14mm}
\begin{minipage}{0.86\linewidth}\small
\textbf{本册特色}\par\smallskip
\begin{itemize}
\item 零基础预备：环境安装、命令行、HTTP 与 SQL 的最小必要知识，术语首现必解释；
\item 标准题面卡：每题写清给定材料、目标、交付物、考点与难度，读完就能动手；
\item 逐题图解带做：W0--W3 每题配有真实终端实录截图与本地靶场真实网页截图，看到截图能照着做；
\item 真实不编造：本地靶场的演示结果来自真实运行；平台题按当前题面和实际反馈记录结论。
\end{itemize}
\end{minipage}
\vspace{12mm}
{\large Web：W0--W3（含题面卡、图解带做与自测）\par}
\vfill
{\large 适用对象：学过 Python、C 和命令行的 CTF 初学者\par}
\vspace{4mm}
{\large 版本：2026 年 9 月\par}
\end{titlepage}

\chapter*{使用说明}
\addcontentsline{toc}{chapter}{使用说明}
本讲义是 Web 方向的专攻讲义：知识点铺垫在前、题目在后，全部内容按零基础读者的接受顺序重新编排。适用对象为学过 Python、C 和命令行、但没有 CTF 经验的初学者。建议先只读第 4 章的题面卡自行尝试，再对照同一题后面的解题步骤和截图复现。

\section*{学习路线图}
\begin{center}
\begin{tabular}{p{0.26\textwidth}p{0.58\textwidth}}
\toprule
章节 & 核心收获 \\
\midrule
第 1 章 零基础预备 & 装好环境、会用命令行、能启动并验证本地靶场、会记证据 \\
第 2 章 Web 基础知识铺垫 & 看懂 HTTP 请求响应、URL 编码、Cookie、SQL 拼接，分清三类漏洞 \\
第 3 章 Web 工具箱 & 会用 PowerShell/WSL、curl.exe、开发者工具、读源码命令 \\
第 4 章 题目：题面卡与逐题图解带做 & 每题一张题面卡，卡后面紧跟解题步骤：照着截图把 W0--W3 逐题做完、做对、说清为什么 \\
第 6 章 复盘与易错清单 & 建立"下次遇到这类题怎么办"的固定流程 \\
第 7 章 自测与参考答案 & 检验概念是否真的掌握 \\
附录 & 平台题目复核、提交 checklist、命令速查表、术语表 \\
\bottomrule
\end{tabular}
\end{center}

\section*{怎么用这本讲义}
\begin{enumerate}
\item \textbf{先做后看}：第 4 章每题先读题面卡，再遮住后面的解题步骤自己动手试，卡住再看“分步操作与观察”；
\item \textbf{照截图核对}：每一步都有"应该看到什么"，输出对不上就看排错指南；
\item \textbf{留下证据}：按 1.5 节的五步记录法保存基线请求和改动后请求，不只记 flag；
\item \textbf{多改一改}：每题末尾的变式练习是提高最快的环节。
\end{enumerate}

\section*{安全声明}
所有实验只针对本资料附带的本地程序（\verb|127.0.0.1|）和题目明确允许练习的授权平台。不要把命令用于未经授权的系统：未经授权的访问测试是违法行为。本讲义中的注入、越权技巧只在本地靶场和授权题目里使用。

\section*{截图说明}
\begin{itemize}
\item \textbf{终端截图}：全部由 \verb|tools/termcap.py| 在真实 PowerShell/WSL 中执行命令、\verb|tools/term2png.py| 渲染生成，每张图在 \verb|figures/raw/| 下留有同名 \verb|.txt| 实录，输出内容没有手打；
\item \textbf{网页截图}：全部是本地靶场 \url{http://127.0.0.1:8765/} 的真实页面，用 Edge 无头模式截取，其中 W1、W2 的页面里能看到真实 flag；
\item \textbf{原理示意图}：由 matplotlib 绘制，画的是结构关系，图注注明"示意图"；
\item \textbf{平台题目}：玄机平台页面不做截图，涉及平台页面事实处使用"示意，非平台截图"的文字表格。
\end{itemize}

\section*{术语约定}
\textbf{flag}：题目的答案字符串，形如 \verb|flag{...}|，提交前检查大小写和花括号。\textbf{payload}：为触发漏洞而构造的输入数据，本讲义译作"注入串"或"载荷"。\textbf{靶场}：专门用来练习的、故意留有漏洞的本地或授权环境。\textbf{参数}：URL 或表单里"名字=值"的一项，是客户端可控的输入。\textbf{基线}：不做任何改动时的正常请求与响应，用来和后面的改动做对比。

\tableofcontents
\mainmatter

% ===========================================================================
\chapter{零基础预备：环境、命令行与本地靶场}
\label{web:prep}
\goal{把环境装好、把命令行用顺、把本地靶场跑起来，并养成"先记录基线"的习惯。本章不讲漏洞，只做准备。}

\section{你需要准备什么}
做本册 Web 题只需要一台 Windows 电脑和 Python 3。按下表准备，右列是"怎么确认装好了"：
\begin{center}
\begin{tabular}{p{0.24\textwidth}p{0.34\textwidth}p{0.32\textwidth}}
\toprule
工具 & 用途 & 怎么确认装好 \\
\midrule
Python 3.8 以上 & 运行本地靶场和演示脚本 & \verb|python --version| 显示版本号 \\
curl & 发送 HTTP 请求、看原始响应 & \verb|curl.exe --version| 显示版本号 \\
Edge 或 Chrome 浏览器 & 看网页、开开发者工具 & 能打开任意本地页面 \\
WSL（可选）& 运行 Linux 命令、做其他方向 & \verb|wsl| 能进入 Ubuntu \\
XeLaTeX（可选）& 编译本讲义 & \verb|xelatex --version| \\
\bottomrule
\end{tabular}
\end{center}

\section{两种命令行：PowerShell 与 WSL}
Windows 自带的命令行叫 \kw{PowerShell}（右键开始菜单就能打开）；WSL 是"Windows 里的 Linux 子系统"。本册 Web 部分全程在 PowerShell 里做即可，WSL 只在需要 Linux 工具时使用。两者的关键差异：
\begin{itemize}
\item 路径写法不同：PowerShell 用 \url{C:/Users/...}，WSL 里同一个位置写成 \url{/mnt/c/Users/...}；
\item 列目录命令不同：PowerShell 用 \verb|dir| 或 \verb|Get-ChildItem|，Linux 用 \verb|ls|；
\item 切换目录都是 \verb|cd|，查看当前目录都是 \verb|pwd|。
\end{itemize}
\figref{web:wsl} 是在真实 WSL 里执行 \verb|uname -a|、\verb|python3 --version|、\verb|which curl| 以及进入资料目录的实录。先看懂它，再决定你要用哪一侧。

\shot{fig-web-13-wsl-basics.png}{WSL 里的真实实录：查看系统、Python、curl 的位置，并进入资料目录（终端实录截图）}{web:wsl}

\pitfall{PowerShell 里不写扩展名就找不到命令。例如直接敲 \code{curl} 会调到 PowerShell 自己的别名（见 3.2 节），要写 \code{curl.exe}；直接敲 \code{python} 一般没问题，但敲 \code{python3} 在 Windows 上可能报"找不到命令"。}

\tip{Windows 终端默认不是 UTF-8，Python 直接打印中文可能乱码。跑本册的 Python 脚本时统一加一个参数 \code{-X utf8} 强制 UTF-8 输出，例如 \code{python -X utf8} \url{figures/raw/fig-web-10-param-fix.py}；第 4 章与附录 B 的命令都按这个写法给。}

\section{启动本地靶场}
本册 W0--W3 的所有练习都在一个本地小服务上做，它就是 \verb|labs/web/app.py|：一个故意留了两个漏洞的 Python 程序，只监听你自己的电脑。在资料根目录执行：

\begin{lstlisting}
python labs/web/app.py
\end{lstlisting}

成功时终端会打印一行 \verb|Lab server: http://127.0.0.1:8765/|（\figref{web:server-start}）。这行的意思是：服务已经在本机的 8765 端口等待请求。术语解释：\kw{服务器}在这里指"等待并回应请求的程序"；\kw{端口}是同一台电脑上区分不同服务的编号；\kw{127.0.0.1} 是"本机自己"的地址（也叫回环地址），别人从网络上访问不到，这是靶场的安全设计。

\shot{fig-web-01-server-start.png}{启动本地靶场：终端打印服务地址，说明已在 127.0.0.1:8765 监听（终端实录截图）}{web:server-start}

启动后按 \verb|Ctrl+C| 停止服务。做完每道题都记得停掉，避免端口被占用。如果启动时报"找不到文件"，先用 \verb|pwd| 看当前目录，再 \verb|cd| 到资料根目录。

\debugtip{服务起不来，按顺序查三件事：① 当前目录是不是资料根目录（\code{pwd} 应该显示 \code{CTF零基础自学资料包} 这个目录）；② Python 是不是装好了（\code{python --version}）；③ 端口是不是被上次没关的服务占了（再跑一次命令，看报错是不是"Address already in use"，是的话关掉旧终端里的服务再试）。}

\section{分诊第一课：先看类型，再看基线}
桌面上给你一份 \verb|labs/misc/pixel.png|，同时给你一个地址 \url{http://127.0.0.1:8765/gate?role=guest}。两条最先该做的检查是什么？答案是：对 PNG 运行 \verb|file labs/misc/pixel.png| 看它到底是什么文件；对网页先启动靶场、再发一次不带任何改动的请求，保存完整响应。前者属于 Misc 方向，第一条证据是文件内容格式；后者属于 Web 方向，第一条证据是 HTTP 请求和响应。

这一步叫\kw{分诊}：判断题目属于哪个方向、该用哪套方法。CTF 常见方向还有 Crypto（密文和数学结构）、Reverse（程序如何判断输入）、Pwn（输入怎样改变内存与控制流）。把网页响应当图片做隐写、见到 PNG 只看画面，都是跳过分诊的典型错误。

\section{五步记录法与证据链}
每次动手都按下面五步走，这是本册所有带做的固定骨架：
\begin{enumerate}
\item \textbf{题面与交付物}：要拿到什么？flag、文件，还是一句结论？
\item \textbf{基线}：不改动输入时，请求和响应长什么样？存下来；
\item \textbf{假设}：我认为哪个输入会被当成"别的东西"（权限、SQL 语法、路径）？
\item \textbf{单变量验证}：一次只改一个输入，观察响应差异；
\item \textbf{复核}：答案能独立验证吗？对源码、重算、或平台反馈。
\end{enumerate}
一条可复现的证据链包括：输入是什么、从哪一行发现异常、用什么命令验证、输出是什么、为什么能推出答案。这样队友能独立复现，你也能分清"猜对了"和"理解了"。

\memory{先类型、后基线、再改一个变量、最后复核——五步走完，题就解完了八成。}

% ===========================================================================
\chapter{Web 基础知识铺垫}
\label{web:basics}
\goal{看懂一次 HTTP 请求与响应的每一行，理解 URL 编码与 Cookie，看懂 SQL 拼接，并能把 SQL 注入、XSS、路径穿越三种问题分开。}

\section{一次 HTTP 请求的解剖}
\kw{HTTP}（超文本传输协议）是浏览器和网站之间"你问我答"的约定：客户端发一个\kw{请求}（request），服务器回一个\kw{响应}（response）。打开网页、提交表单、App 刷数据，底层都是 HTTP 请求。\figref{web:http-anatomy} 把本地靶场的一次真实交换逐段拆开（示意图，内容取自实测）。

\shot{fig-web-19-http-anatomy.png}{一次 HTTP 请求的解剖图：上为请求、下为响应，右侧逐段解释（原理示意图，内容与本地靶场实测一致）}{web:http-anatomy}

请求分四段：\kw{请求行}（方法 + 路径 + 查询串 + 协议版本）、\kw{请求头}（每行"名字: 值"）、空行、\kw{请求体}。响应分四段：\kw{状态行}（协议 + 状态码 + 短语）、\kw{响应头}、空行、\kw{响应体}。响应体就是网页内容，flag 通常藏在这里。空行是硬性规定：它把"头"和"体"分开，请求和响应都遵守。

响应头里有几个反复出现的字段，认识它们能帮你判断"哪里不对劲"：\verb|Content-Type| 说明响应体是什么格式（\verb|text/html| 是网页、\verb|text/plain| 是纯文本、\verb|application/json| 是数据接口），\verb|Content-Length| 说明正文有多少字节，\verb|Location| 用在跳转（301/302）时给出新地址。

方法（method）常见的有：
\begin{itemize}
\item \kw{GET}：取数据，参数通常放在 URL 的问号后面，例如 \verb|/gate?role=guest|；
\item \kw{POST}：交数据，参数通常放在请求体里，例如提交上传表单；
\item 其他还有 PUT、DELETE、HEAD 等，本册只需要认识前两个。
\end{itemize}

状态码（status code）是响应开头的三位数字，做题时先看它再看正文：
\begin{center}
\begin{tabular}{llp{0.52\textwidth}}
\toprule
状态码 & 含义 & 典型场景 \\
\midrule
200 & 成功 & 正常返回页面或数据 \\
301/302 & 跳转 & 响应头 \verb|Location| 指向新地址 \\
400 & 请求有问题 & 服务器认为你发的格式不对 \\
401/403 & 未认证 / 无权限 & 登录失效、权限不足 \\
404 & 找不到 & 路径写错，或该页面不存在 \\
500 & 服务器出错 & 程序抛异常，报错信息可能泄露线索 \\
\bottomrule
\end{tabular}
\end{center}

\pitfall{浏览器地址栏只展示了请求的一小部分。真实请求还有请求头、Cookie、请求体。用 F12 或 curl 看到的才是全貌。}

\memory{请求=方法+路径+查询串+头+体；响应=状态码+头+体；中间的空行把头和体分开。}

\section{URL 与 URL 编码}
一个 URL 的结构是：\verb|协议://主机:端口/路径?查询串#片段|。以 \url{http://127.0.0.1:8765/gate?role=guest} 为例：协议 \verb|http|，主机和端口 \verb|127.0.0.1:8765|，路径 \verb|/gate|，查询串 \verb|role=guest|。查询串里可以有多个参数，用 \verb|&| 连接，例如 \verb|a=1&b=2|。

URL 里有些字符有特殊含义（空格、问号、井号、引号……），想把它们当普通数据传，就要做\kw{URL 编码}：把字节写成 \verb|%XX| 的形式。常见对照：
\begin{center}
\begin{tabular}{lll}
\toprule
字符 & URL 编码 & 说明 \\
\midrule
空格 & \verb|%20| & 也常写成 \verb|+|（表单编码里） \\
单引号 & \verb|%27| & SQL 注入里最关键的一个 \\
等号 & \verb|%3D| & 条件 \verb|1=1| 里的等号 \\
斜杠 & \verb|%2F| & 路径相关题目会遇到 \\
井号 & \verb|%23| & URL 片段分隔符，不编码会被浏览器截断 \\
\bottomrule
\end{tabular}
\end{center}

\pitfall{URL 编码既不是加密，也不是防护。服务器收到后会解码还原成原字符。把单引号写成 \code{\%27}，数据库看到的仍然是单引号。编码只解决"传输"问题，不解决"信任"问题。}

\memory{编码只管运输，不管安全；\code{\%27} 到了服务器还是单引号。}

\section{Cookie 与会话}
服务器怎么知道"这个请求和上一个请求是同一个人"？HTTP 本身不记人，靠 \kw{Cookie} 补上：服务器在响应头发 \verb|Set-Cookie|，浏览器存下来，之后每个请求都自动带上。Cookie 里常放一个\kw{会话标识}（session id），服务器用它找到"你登录过了"的记录。

但要记住：Cookie 存在浏览器里，值由客户端送来，客户端可以改它、可以伪造请求带上任意 Cookie。所以"请求里带着 Cookie"不等于"请求可信"。类似地，隐藏表单域、按钮置灰、把参数从 URL 挪到 POST，都不改变"客户端可控"这个事实。

\memory{凡是客户端送来的（URL 参数、POST 表单、Cookie、请求头）都是用户填的，服务端都要自己验证。}

\section{SQL 基础：语句结构与字符串拼接}
\kw{SQL} 是操作数据库的语言。把数据库想成一组表格：\kw{表}有名字，表里一行是一条记录，每列是一个字段。本册靶场有一张 \verb|notes| 表，两列：\verb|title|（标题）和 \verb|body|（内容），共两行：\verb|public| 和 \verb|staff|。

查询一行记录的语句长这样：
\begin{lstlisting}
SELECT title, body FROM notes WHERE title = 'public'
\end{lstlisting}
\kw{SELECT} 说"要哪几列"，\kw{FROM} 说"从哪张表"，\kw{WHERE} 说"要满足什么条件"。字符串用单引号包围：\verb|'public'| 是一个完整的字符串值。

\kw{WHERE} 后面是一个\kw{条件表达式}，可以先比大小、再用逻辑词组合：
\begin{itemize}
\item 比较：\verb|=|（相等）、\verb|!=| 或 \verb|<>|（不等）、\verb|>|、\verb|<|；
\item 逻辑：\verb|AND|（两边都真，结果才真）、\verb|OR|（有一边真，结果就真）、\verb|NOT|（取反）。
\end{itemize}
例如 \verb|WHERE title = 'public'| 只匹配标题正好是 public 的那一行；而 \verb|WHERE 1=1| 永远为真（1 恒等于 1），会匹配\emph{所有}行；\verb|WHERE 1=2| 永远为假，一行都不匹配。这两句本身毫无意义，却是注入里最常借用的"恒真 / 恒假"条件——第 4 章 W2 的 \verb|OR 1=1| 用的就是这个。另外，SQL 关键字\emph{大小写不敏感}，\verb|select| 和 \verb|SELECT| 完全等价。

SQL 还有两个会经常打照面的写法，这里先混个眼熟。一是 \kw{LIMIT}，限定返回几行，如 \verb|SELECT title FROM notes LIMIT 1| 只取第一条；二是 \kw{UNION SELECT}，把另一条查询的结果上下拼接到当前结果后面，如 \code{SELECT title, body FROM notes UNION SELECT 'a', 'b'}。用 \verb|UNION| 有一条硬规矩：\textbf{两条查询的列数必须相同}，否则数据库会直接报"列数不匹配"的错误。本册正文不靠它们解题，但读网上题解时会遇到，别被吓到。

\kw{注释}是"让数据库忽略掉的一段文字"：从注释符开始到行尾，数据库一律不看。SQL 里常见两种注释符：
\begin{itemize}
\item \verb|--|：两个连字符，后面必须紧跟一个空格（或换行）才算注释，稳妥写法是本文的 \verb|-- |（后面留一个空格），本册注入串末尾用的就是它；
\item \verb|#|：只在 MySQL 这类部分数据库里是注释，SQLite 不认它，换到别的数据库题里别想当然。
\end{itemize}
注释在注入里的用处是"把原语句尾部多出来的内容吃掉"，让拼出来的语句重新变回合法语法。

\pitfall{注释符由两个连字符组成，而且后面必须跟一个空格（或换行）。写成紧贴文字的 \code{--x} 时，有的数据库不把它当注释，注入就会失败。记牢：\code{--} 之后再敲一个空格。}

靶场 \verb|labs/web/app.py| 第 39 行用 f-string 把用户输入拼进这条语句：
\begin{lstlisting}[language=Python]
sql = f"SELECT title, body FROM notes WHERE title = '{title}'"
\end{lstlisting}
如果 \verb|title| 是 \verb|public|，拼出来一切正常。如果 \verb|title| 是一个单引号，拼出来的语句就多了个没配对的引号，SQL 语法被破坏，数据库报错。这就是\kw{字符串拼接}的危险：用户数据被拼进了语句\emph{结构}，而不只是语句里的一个\emph{值}。

正确做法叫\kw{参数化查询}（也叫预编译占位符）：语句结构先定死，数据用占位符 \verb|?| 隔开，单独提交：
\begin{lstlisting}[language=Python]
db.execute("SELECT title, body FROM notes WHERE title = ?", (title,))
\end{lstlisting}
此时无论 \verb|title| 里写什么，数据库都把它当作一个完整的标题值。

把这两种写法放进同一个内存数据库跑一遍：同一张表、同一个拼接位置，结果完全不同（\figref{web:sql-normal} 与 \figref{web:sql-inject}，均为真实执行，只操作内存数据库、不启动靶场，所以不会占用端口）：

\begin{lstlisting}
python -X utf8 figures/raw/fig-web-20-sqlite-demo.py normal
python -X utf8 figures/raw/fig-web-20-sqlite-demo.py inject
\end{lstlisting}

\shot{fig-web-20-sql-normal.png}{正常查询：title=public，拼出的语句合法，只返回公开的那一条（真实执行）}{web:sql-normal}

\shot{fig-web-21-sql-inject.png}{注入输入：同一个拼接位置，输入里的单引号提前收尾、OR 1=1 变成恒真条件、注释符吃掉尾部引号，两行全返回（真实执行）}{web:sql-inject}

两张图对比着看"拼出来的 SQL"那一行就很清楚：正常输入时，数据老实待在两个引号之间；注入输入时，数据自带的引号跑到了引号\emph{外面}，把 \verb|OR 1=1| 顶成了语句的一部分。这就是"拼接把数据变成了结构"最直观的样子，也正是参数化查询要解决的问题。

\memory{拼接=数据混进语法；参数化=语法和数据分家。SQL 注入的全部原理就这一句。}

\section{SQL 注入、XSS、路径穿越：三种"数据被当成代码"}
初学最常犯的错是见到输入点就喊"SQL 注入"。三者的区别在于\textbf{数据被谁解释成了代码}：
\begin{center}
\begin{tabular}{p{0.14\textwidth}p{0.24\textwidth}p{0.24\textwidth}p{0.26\textwidth}}
\toprule
类型 & 被哪个组件解释 & 典型现象 & 典型修复 \\
\midrule
SQL 注入 & 数据库 & 输入单引号报 SQL 错；条件恒真多返回记录 & 参数化查询 \\
XSS（跨站脚本）& 浏览器 & 输入的脚本在别人的页面上执行 & 按输出位置做 HTML 转义 \\
路径穿越 & 文件系统 & 读到目录外的文件（如 \verb|../etc/passwd|）& 规范化路径后核对根目录 \\
\bottomrule
\end{tabular}
\end{center}

\textbf{XSS}（跨站脚本）：应用把用户输入原样拼进网页，浏览器把这段数据当成 HTML 或脚本。它发生在"浏览器解释输出"的阶段，与 SQL 注入发生在"数据库解释查询"的阶段不同。修复要按输出位置（HTML 文本、属性、脚本）做编码，仅在输入端"禁止尖括号"通常会被绕过。

\textbf{路径穿越}：程序把用户给的文件名拼到某个目录后面就打开文件，用户用 \verb|../| 走出允许的目录。修复不是删掉某个固定字符串，而是把路径规范化后核对它仍位于允许的根目录内，并在访问控制层确认用户有权读它。

判断方法很朴素：看这个输入最终被谁消费。进了 SQL 语句 → SQL 注入；进了网页输出 → XSS；进了文件路径 → 路径穿越；进了权限判断 → 越权（下一节）。

\pitfall{同一个参数可能同时通向多个解释器。"输入被过滤过"不等于安全：过滤清单可以被编码、大小写、双写等方式绕过。可靠的修复是让解释器收到的数据永远只是数据。}

\section{认证、授权与会话管理}
网站至少要区分三个问题：
\begin{itemize}
\item \textbf{认证}（你是谁）：核对账号口令或令牌，确认身份；
\item \textbf{授权}（你能做什么）：根据身份和资源归属，判断这次能不能做；
\item \textbf{会话管理}（这次请求属于哪次登录）：把一系列请求和同一身份绑定。
\end{itemize}
W1 的问题不在"URL 上出现了 role"，而在于服务器把客户端自报的 \verb|admin| 当成了授权结论。隐藏按钮、禁用控件、改用 POST，都只是界面层的限制；只要服务器的路由还接受这个值，越权就在。真实系统还要核对资源归属：用户 A 登录了也不该看到用户 B 的私有数据（这类问题常叫"水平越权"或"越权访问"，练习题里常见线索是可改的对象编号）。

分析权限题时画一张小表：匿名用户、普通用户、管理员分别能访问哪些资源；哪些判断在页面上做（不可信），哪些在服务端路由里做（才算数）。

\memory{认证回答"你是谁"，授权回答"你能做什么"；两者都在服务端完成才有意义。}

% ===========================================================================
\chapter{Web 工具箱}
\label{web:tools}
\goal{每个工具知道是什么、怎么调用、输出怎么读，并认识 Windows 上的两个大坑（curl 别名、PowerShell 引号）。}

\section{PowerShell 与 WSL 基本命令}
在 PowerShell 里做 Web 题，常用命令就这几个（WSL 对照见 \figref{web:wsl}）：
\begin{center}
\begin{tabular}{p{0.3\textwidth}p{0.3\textwidth}p{0.3\textwidth}}
\toprule
想做什么 & PowerShell & WSL/Linux \\
\midrule
看当前目录 & \verb|pwd| & \verb|pwd| \\
列文件 & \verb|dir| & \verb|ls| \\
切换目录 & \verb|cd 路径| & \verb|cd 路径| \\
看文件内容 & \verb|Get-Content 文件| & \verb|cat 文件| \\
搜文件内容 & \verb|Select-String| & \verb|grep| \\
清屏 & \verb|cls| & \verb|clear| \\
\bottomrule
\end{tabular}
\end{center}
路径里含空格时整个路径加引号。命令敲错时先看报错第一行，Windows 的报错很长，真正的信息通常在开头。

\section{curl：把请求原样发出去}
\kw{curl} 是命令行的 HTTP 客户端：你告诉它 URL 和选项，它把请求发出去并把响应打印出来。Web 题里它是最重要的工具，因为浏览器"帮忙"做的太多（自动编码、自动带 Cookie、自动跟随跳转），而 curl 让你看清每个字节。常用选项：
\begin{center}
\begin{tabular}{p{0.18\textwidth}p{0.68\textwidth}}
\toprule
选项 & 作用 \\
\midrule
\verb|-s| & 静默模式，去掉进度条噪音 \\
\verb|-i| & 打印响应头（状态行 + 响应头 + 正文）一起看 \\
\verb|-G| & 把 \verb|--data| 的内容拼到 URL 查询串上（发 GET） \\
\code{--data-\allowbreak urlencode} \code{k=v} & 对 v 做 URL 编码后再发送，避免手工编码出错 \\
\verb|-X POST| & 指定方法为 POST \\
\verb|-d "a=1"| & 放入请求体（常与 POST 连用） \\
\code{-H} \code{"Name: value"} & 添加一个请求头 \\
\verb|-v| & 显示完整交互过程（含发出的请求头） \\
\bottomrule
\end{tabular}
\end{center}

\pitfall{Windows 的 PowerShell 里 \code{curl} 不是你以为的那个 curl！它是 \code{Invoke-WebRequest} 的别名。要调真正的 curl，必须写 \code{curl.exe}（带扩展名）。\figref{web:curl-alias} 是真实验证：\code{Get-Alias curl} 显示它指向 \code{Invoke-WebRequest}，而 \code{curl.exe --version} 才显示 libcurl 版本号。}

\shot{fig-web-11-curl-vs-alias.png}{Windows 实录：curl 是 Invoke-WebRequest 的别名，真正 curl 要写 curl.exe（终端实录截图）}{web:curl-alias}

还有一个 Windows 专属的坑：PowerShell 对引号的处理和 bash 不同，直接写
\verb|--data-urlencode "title=' OR 1=1 -- "| 时内层引号可能被提前吃掉。稳妥写法是用单引号字符串，字符串内部的单引号写两个：
\begin{lstlisting}
curl.exe -s -i -G --data-urlencode 'title='' OR 1=1 -- ' http://127.0.0.1:8765/notes
\end{lstlisting}
看到"字符串缺少终止符"这类报错，说明是 PowerShell 解析出错，命令还没发出去——这不是 SQL 的错。第 4 章截图里的命令就是这个写法。

\section{浏览器开发者工具（概念）}
按 \verb|F12| 打开开发者工具，Web 题最常用 \kw{Network}（网络）面板。它能告诉你：
\begin{itemize}
\item 每个请求的方法、完整 URL（含查询串）、状态码；
\item 请求头、Cookie、请求体（点开单条请求看）；
\item 响应头和响应体（和 curl 看到的一致）；
\item 页面加载时到底发了几个请求（静态文件、接口、统计脚本各算一个）。
\end{itemize}
用好它的关键是"和 curl 互相印证"：在面板里看到可疑请求后，右键可以复制成 cURL 命令，粘到终端里重放、改参数。开发者工具不能修改服务器的判断——页面上按钮变灰只是界面效果，理解这一点就不会把"看不见"当成"不存在"。

\section{\texttt{python -m http.server} 与靶场的区别}
有人问：直接用 \verb|python -m http.server| 不就能开网站吗？能，但它只是\kw{静态文件服务}：把目录里的文件原样发出去，没有任何逻辑。\figref{web:http-server} 的实录显示它对 \verb|app.py| 返回的是文件源码本身。而本地靶场 \verb|app.py| 是\kw{动态应用}：它按路径分发（\verb|/gate|、\verb|/notes|）、读参数、查数据库、按条件返回不同内容——漏洞就藏在这些"按参数做决定"的逻辑里。做 Web 题时先问自己：这个站点是"发文件"还是"做判断"？

\shot{fig-web-12-http-server.png}{python -m http.server 实录：它只是把文件原样发出（返回 app.py 的源码），没有动态判断（终端实录截图）}{web:http-server}

\section{读源码：Get-Content 与 Select-String}
Web 题常常附源码，或者干脆让你自己审计本地文件。PowerShell 里两个命令足够：
\begin{lstlisting}
Get-Content labs/web/app.py                      # 看整个文件
Select-String -Path labs/web/app.py -Pattern 'sql = f' -Context 1,2   # 搜关键行并显示上下文
\end{lstlisting}
找漏洞时先搜"数据被拼接的地方"：搜 \verb|SELECT|、搜 \verb|f"|、搜 \verb|execute|、搜 \verb|open(|。命中后向前后各看几行，问一句："这个变量的值是哪来的？"——从 \verb|query| / \verb|request| / \verb|Cookie| 来的，就是外部输入。\figref{web:app-source} 就是这么找到靶场里那一行的。

\shot{fig-web-09-app-source.png}{阅读 labs/web/app.py：/notes 的处理逻辑，以及被 Select-String 高亮的拼接行（终端实录截图）}{web:app-source}

\memory{浏览器看现象，curl 看字节，源码看原因——三样凑齐，漏洞就解释得清。}

% ===========================================================================
% ===========================================================================
\chapter{题目：题面卡与逐题图解带做}
\label{web:cards}
\goal{四道题，每题一节：先用题面卡把“给什么、要什么、考什么”读清楚，紧接着就是这道题的完整解题步骤——读题与思路、分步操作与观察、结果与核对、为什么会成功、排错指南、变式练习。建议先只读题面卡自己动手试 10--15 分钟，卡住再看后面的步骤。}

每张卡的字段固定为：题目来源 / 给定材料 / 目标 / 交付物 / 考点 / 前置知识 / 难度。难度用星表示：$\bigstar\circ\circ$ 入门、$\bigstar\bigstar\circ$ 进阶、$\bigstar\bigstar\bigstar$ 挑战。四道题都只用本册靶场，离线可做。

% ===========================================================================
\section{题 W0：谁决定了 role 的值？}

\begin{challenge}{题 W0 题面卡}
\field{题目来源}{本地靶场 \code{labs/web/app.py}（离线可做）}
\field{给定材料}{一段真实请求文本：\code{GET /gate?role=guest HTTP/1.1} 与 \code{Host: 127.0.0.1:8765}；运行中的本地服务}
\field{目标}{把请求拆成方法、路径、参数名、参数值四部分，圈出由客户端控制的部分，并保存一份不改动的基线响应}
\field{交付物}{一张标注好的请求拆解 + 一份基线响应记录（本题不交 flag）}
\field{考点}{HTTP 请求结构；"客户端可控"意识}
\field{前置知识}{2.1 节 HTTP 请求解剖、3.2 节 curl、1.3 节启动靶场}
\field{难度}{$\bigstar\circ\circ$（入门，纯观察无 flag）}
\end{challenge}

\paragraph{读题与思路}
题面给了一段请求文本。先别急着找漏洞——这题练的是"看清客户端能控制什么"。看请求时按顺序问自己四个问题：用什么\textbf{方法}？要什么\textbf{路径}？带了哪些\textbf{参数}？哪个部分是我可以改的？浏览器地址栏显示的只是请求的一小部分，完整请求要看 curl 或开发者工具。

\paragraph{分步操作与观察}
\begin{enumerate}[label=\arabic*.]
\item 启动本地靶场（1.3 节的命令），确认终端打印服务地址（\figref{web:server-start}）。
\item 用浏览器打开 \url{http://127.0.0.1:8765/}，看到一行提示："Try /gate?role=guest and /notes?title=public"（\figref{web:home-browser}）。这是靶场的"入口索引"，以后打开陌生站点也先找这种索引页。
\item 在 PowerShell 里执行下面的命令，保存基线响应，应该看到状态行 \verb|HTTP/1.0 200 OK| 和正文提示（\figref{web:home-curl}）：
\begin{lstlisting}
curl.exe -s -i http://127.0.0.1:8765/
\end{lstlisting}
\item 对照 \figref{web:http-anatomy} 拆解请求：方法是 \verb|GET|，路径是 \verb|/|，没有查询参数；响应里有状态行、\verb|Content-Type| 等响应头、空行、正文。
\item 访问门禁接口，把响应存下来，正文是 \verb|Access denied|（\figref{web:gate-guest-curl}）：
\begin{lstlisting}
curl.exe -s -i http://127.0.0.1:8765/gate?role=guest
\end{lstlisting}
浏览器里的同一页见 \figref{web:gate-guest}。
\item 在请求文本上标注：方法 \verb|GET|、路径 \verb|/gate|、参数名 \verb|role|、参数值 \verb|guest|。其中\textbf{参数值由发请求的一方决定}，这就是"客户端可控"的部分。
\end{enumerate}

\shot{fig-web-14-home-browser-crop.png}{靶场首页（真实网页截图）：浏览器访问 http://127.0.0.1:8765/ 的实际页面（裁切显示内容）}{web:home-browser}

\shot{fig-web-02-home-curl.png}{基线请求：curl.exe -s -i 访问首页，打印状态行、响应头与正文（终端实录截图）}{web:home-curl}

\shot{fig-web-03-gate-guest-curl.png}{请求 /gate?role=guest：正文是 Access denied，这就是 W1 要对比的基线（终端实录截图）}{web:gate-guest-curl}

\shot{fig-web-15-gate-guest-crop.png}{同一页面的浏览器视图（真实网页截图）：/gate?role=guest 返回 Access denied（裁切显示响应文本）}{web:gate-guest}

\paragraph{结果与核对}
本题交付物是"标注好的请求拆解 + 基线记录"，没有 flag。核对方法：把你的四部分标注和上面第 6 步对照；基线文件里应包含状态码 \verb|200|、响应头和完整正文。基线存在的意义是：以后每次改动输入，都有一个可以对比的"正常样子"。

\paragraph{为什么会成功}
因为 HTTP 请求的每个字段都是客户端构造的。浏览器只是"帮你发请求的程序"，它不提供任何可信度；用 curl、脚本、抓包工具都能重建这个 GET 请求。所以读 Web 题的第一课就是：分清哪些字段是服务端定的（路径对应的功能），哪些是客户端填的（参数、头、Cookie）。

\debugtip{① 连接被拒绝（curl 报 Connection refused）：靶场没启动，或不在 8765 端口。② 状态码 404：路径敲错，注意大小写和开头的斜杠。③ 中文乱码：PowerShell 编码问题，命令前先执行 \code{[Console]::OutputEncoding=[System.Text.Encoding]::UTF8}。④ 输出一大堆进度条：忘了加 \code{-s}。}

\paragraph{变式练习}
\begin{enumerate}
\item 把参数值改成 \verb|administrator|，记录响应是否变化——"响应变了"说明服务端读了这个参数，但还不能断言有漏洞。
\item 把 \verb|?role=guest| 整段删掉再请求，观察服务端的默认行为（靶场会按缺省值处理）。
\item 用 \verb|curl.exe -v| 重发第 5 步的请求，对照发出的请求头和图解中的"请求头"部分。
\end{enumerate}

% ===========================================================================
\section{题 W1：谁说你是管理员？}

\begin{challenge}{题 W1 题面卡}
\field{题目来源}{本地靶场 \code{labs/web/app.py} 的 \code{/gate} 接口（离线可做）}
\field{给定材料}{接口 \code{/gate?role=guest} 只返回普通内容；目标是读出管理员内容}
\field{目标}{只修改一个输入，让接口返回管理员内容，并解释这个输入为什么生效}
\field{交付物}{一个 \code{flag\{...\}} 字符串，提交前检查大小写与花括号；外加一句话漏洞解释}
\field{考点}{授权逻辑：服务端信任了客户端自报的身份字段}
\field{前置知识}{2.6 节认证与授权、2.1 节查询参数、3.2 节 curl}
\field{难度}{$\bigstar\circ\circ$（入门）}
\end{challenge}

\paragraph{读题与思路}
目标是读出"管理员内容"，而入口只有一个带参数的接口。按五步法：基线已经有了（\verb|Access denied|），假设是"服务器把 \verb|role| 参数当成权限结论"，验证方法是\textbf{只改参数值}，其他一律不动。一次只改一个变量，才能说清是哪个输入造成的变化。

\paragraph{分步操作与观察}
\begin{enumerate}[label=\arabic*.]
\item 先请求 guest 留基线（\figref{web:gate-guest-curl}），记下正文 \verb|Access denied|。
\item 只把参数值从 \verb|guest| 改成 \verb|admin|，其余不动，执行：
\begin{lstlisting}
curl.exe -s -i http://127.0.0.1:8765/gate?role=admin
\end{lstlisting}
响应正文变为 \verb|flag{http_query_is_input}|（\figref{web:gate-admin-curl}）。
\item 浏览器里再看一次同一页（\figref{web:gate-admin}），确认这不是终端里手工打印的文本，而是服务器返回的。
\item 打开源码核对原因，看 \verb|/gate| 分支（\figref{web:app-source}）：\verb|role| 的值直接与 \verb|"admin"| 比较，成立就输出 flag。没有任何登录、口令或会话检查。
\end{enumerate}

\shot{fig-web-04-gate-admin-curl.png}{只改一个输入：role=admin，响应体给出 flag（终端实录截图）}{web:gate-admin-curl}

\shot{fig-web-16-gate-admin-crop.png}{浏览器里的同一请求（真实网页截图）：页面上就是管理员内容 flag（真实截图，裁切显示响应文本）}{web:gate-admin}

\paragraph{结果与核对}
\flagbox{flag\{http\_query\_is\_input\}}
独立核对：把 flag 与 \verb|labs/web/app.py| 中 \verb|/gate| 分支的返回值逐字符对照（见 \figref{web:app-source} 里的高亮行）。flag 只在本地靶场出现，别拿到别的平台提交。

\paragraph{为什么会成功}
机制一句话：\textbf{授权判断读取了客户端可控的字段}。展开说：服务器收到请求后，本应根据"已认证的会话"决定角色，却直接取了查询参数 \verb|role| 的值当结论。因为请求由客户端构造，任何人都能写 \verb|role=admin|。这类漏洞叫\kw{越权}（broken access control），在真实网站里通常表现为改用户编号看别人资料、改角色字段进管理后台。

\debugtip{① 返回 \code{Access denied}：确认 URL 是 \code{/gate?role=admin}，注意参数名是 role、值是 admin（全小写）。② 404：路径写成了 \code{/Gate} 或多了空格。③ flag 显示乱码或缺花括号：是终端显示问题，把响应存文件再看（\code{curl.exe -s URL > out.txt}）。④ 怎么确认不是自己打错字：重跑一遍命令，让服务器"再给一次"。}

\paragraph{变式练习}
\begin{enumerate}
\item 把参数搬到请求体（\verb|-X POST -d "role=admin"|）再试。注意本靶场只读查询参数，所以会失败——但真实网站若把角色放 POST/隐藏域，同样是漏洞。想一想为什么。
\item 用 \verb|-H "Cookie: role=admin"| 把角色放进 Cookie 再请求 guest，观察结果。结论：改传输位置不修复信任问题；本靶场甚至不读 Cookie。
\item 自己出题：写一个按请求里 \verb|is_admin=true| 决定是否展示测试内容的小服务，再想修复方案（提示：服务端会话角色）。
\end{enumerate}

% ===========================================================================
\section{题 W2：标题搜索里，哪个字节改变了查询？}

\begin{challenge}{题 W2 题面卡}
\field{题目来源}{本地靶场 \code{labs/web/app.py} 的 \code{/notes} 接口（离线可做）}
\field{给定材料}{接口 \code{/notes?title=...}；源码 \code{labs/web/app.py} 中拼 SQL 的一行}
\field{目标}{让接口返回本来查不到的 staff 记录，读出其中的 flag；并说明输入的哪一部分被数据库当成了语法}
\field{交付物}{一个 \code{flag\{...\}} 字符串 + 一句"哪段输入进入了 SQL 结构"的解释}
\field{考点}{SQL 注入：外部数据进入查询结构}
\field{前置知识}{2.4 节 SQL 语法结构、2.2 节 URL 编码、3.2 节 curl 的编码选项}
\field{难度}{$\bigstar\circ\circ$（入门）}
\end{challenge}

\paragraph{读题与思路}
这题的线索是"报错随输入改变"。顺序永远是：正常输入建基线 → 单引号探测 → 读源码确认输入的去向 → 手工代入语句推结构 → 构造输入。报错本身不是 flag，但它告诉你：输入进入了某个"会做语法分析"的地方——这里是 SQL。

\paragraph{分步操作与观察}
\begin{enumerate}[label=\arabic*.]
\item 建基线：查正常标题 \verb|public|，只应返回一条公开笔记（\figref{web:notes-baseline}）；浏览器视图见 \figref{web:notes-public}。
\begin{lstlisting}
curl.exe -s -i -G --data-urlencode "title=public" http://127.0.0.1:8765/notes
\end{lstlisting}
\item 探测：发一个单引号，响应变成 SQL 报错；再试 \verb|public'| 也报错（\figref{web:quote-error}）。说明这个引号被数据库读进了语句结构。
\begin{lstlisting}
curl.exe -s -G --data-urlencode 'title=''' http://127.0.0.1:8765/notes
\end{lstlisting}
\item 读源码（\figref{web:app-source}）：关键行是
\begin{lstlisting}[language=Python]
sql = f"SELECT title, body FROM notes WHERE title = '{title}'"
\end{lstlisting}
把输入代进去手写一遍：\verb|title| 填 \verb|public| 得 \verb|WHERE title = 'public'|；填单引号得 \verb|WHERE title = '''|，引号没配对，所以报错。
\item 构造输入 \verb|' OR 1=1 -- |：首引号结束字符串，\verb|OR 1=1| 是恒真条件，\verb|--| 把语句尾部剩下的 \verb|'| 变成注释。发送（\figref{web:sqli-curl}）：
\begin{lstlisting}
curl.exe -s -i -G --data-urlencode 'title='' OR 1=1 -- ' http://127.0.0.1:8765/notes
\end{lstlisting}
\item 浏览器里用 URL 编码形式访问同一输入（\verb|%27| 是单引号、\verb|%20| 是空格、\verb|%3D| 是等号），看到同样两行输出（\figref{web:sqli-urlencoded}、\figref{web:notes-sqli}）：
\begin{lstlisting}
http://127.0.0.1:8765/notes?title=%27%20OR%201%3D1%20--%20
\end{lstlisting}
\end{enumerate}

\shot{fig-web-05-notes-baseline.png}{基线：正常标题 public 只返回一条公开笔记（终端实录截图）}{web:notes-baseline}

\shot{fig-web-17-notes-public-crop.png}{同一基线的浏览器视图（真实网页截图）：/notes?title=public（裁切显示响应文本）}{web:notes-public}

\shot{fig-web-06-notes-quote-error.png}{单引号探测：响应从"正常笔记"变成 SQL 报错，报错随输入改变（终端实录截图）}{web:quote-error}

\shot{fig-web-07-notes-sqli-curl.png}{注入成功：两行记录都返回，包含 staff 的 flag（终端实录截图）}{web:sqli-curl}

\shot{fig-web-08-notes-sqli-urlencoded.png}{同一注入串的 URL 编码形式：\%27、\%20、\%3D，结果一致（终端实录截图）}{web:sqli-urlencoded}

\shot{fig-web-18-notes-sqli-crop.png}{浏览器里的注入结果（真实网页截图）：页面显示两条笔记与 flag（裁切显示响应文本）}{web:notes-sqli}

\paragraph{结果与核对}
\flagbox{flag\{sql\_parameters\_matter\}}
独立核对：这条记录来自服务器数据库的 \verb|staff| 行（建表时插入，见 \verb|app.py| 的 \verb|make_db|），不是终端打印的文本——换 curl、换浏览器都能拿到同一个值，这就是独立复现。

\paragraph{为什么会成功}
机制一句话：\textbf{外部数据进入了 SQL 的语句结构}。展开说：程序用字符串拼接构造查询，输入里的单引号提前结束了字符串字面量，后面的 \verb|OR 1=1| 成为 SQL 的条件语法，\verb|--| 是 SQL 的注释符号，把原本尾部的引号注释掉，语句仍然合法。数据库于是把"标题为空或恒真"当作条件，两行全返回。注意：注入是否成功取决于原语句的引号、注释写法和数据库方言，记住这一条固定输入解不了所有 SQL 题。

\debugtip{① 看到"字符串缺少终止符 / ParserError"：这是 PowerShell 解析错误，命令没发出去。把参数写成单引号字符串，内部单引号写两个（见 3.2 节）。② 返回 \code{No note found}：注入串被当成普通标题了，检查引号是否真的发送了（用 \code{--data-urlencode} 或看 URL 里的 \code{\%27}）。③ 报错和图里不一样：确认查的是 \code{/notes} 而不是 \code{/gate}。④ 想确认服务器收到的到底是什么：把注入串换成 \code{AAAA}，看响应正文是否搜索 AAAA 标题。}

\paragraph{变式练习}
\begin{enumerate}
\item 不用注释符，改用 \verb|public' OR '1'='1| 试试：想一想为什么这个写法不需要 \verb|--|（提示：尾部引号变成了第二个字符串的开头）。
\item 只发一个单引号的 URL 编码形式 \verb|%27|，对照 \figref{web:quote-error} 确认现象一致。
\item 把条件改成 \verb|' AND 1=2 -- |，预测返回几行（提示：恒假条件），再实际验证。
\end{enumerate}

% ===========================================================================
\section{题 W3：同样的输入为什么在修复版里失效？}

\begin{challenge}{题 W3 题面卡}
\field{题目来源}{本地靶场 \code{labs/web/app.py} 的修复实验（离线可做）}
\field{给定材料}{修复写法：\code{db.execute("SELECT ... WHERE title = ?", (title,))}；W2 的注入串}
\field{目标}{把 W2 的查询改成参数化形式，重新发送同一个注入串，预测并验证返回多少条笔记；再确认正常标题 public 不受影响}
\field{交付物}{一组对照结果：恶意输入失效 + 正常输入仍可用，并说明原因}
\field{考点}{参数化查询：让数据永远只是数据}
\field{前置知识}{2.4 节参数化查询、4.3 节 W2 带做}
\field{难度}{$\bigstar\circ\circ$（入门）}
\end{challenge}

\paragraph{读题与思路}
修复不是"删掉危险词"，而是让数据库永远把输入当数据。预测先于操作：注入串在参数化写法里应该只是"一个完整标题"，匹配不到任何记录；正常标题 \verb|public| 应该照常工作。验证修复必须做\textbf{两组测试}：恶意输入失效 + 正常功能保留，缺一不可。

\paragraph{分步操作与观察}
\begin{enumerate}[label=\arabic*.]
\item 把查询改为参数化形式（这就是修复）：
\begin{lstlisting}[language=Python]
db.execute("SELECT title, body FROM notes WHERE title = ?", (title,))
\end{lstlisting}
\item \textbf{重启服务}：不重启的话，运行的还是旧代码。停掉旧进程，再执行 \code{python} \url{labs/web/app.py}。
\item 查询 \verb|public|，确认公开笔记仍正常返回。
\item 用与 W2 \textbf{完全相同}的注入串再查一次，确认只返回 0 行（找不到），staff 记录不再出现。
\item 懒得改服务器也没关系：配套演示脚本把两种写法放在同一个进程里对照跑，实录见 \figref{web:param-fix}（脚本在 \verb|figures/raw/fig-web-10-param-fix.py|，只操作内存数据库）：
\begin{lstlisting}
python -X utf8 figures/raw/fig-web-10-param-fix.py
\end{lstlisting}
\end{enumerate}

\shot{fig-web-10-param-fix.png}{修复前后对照：同一注入串，拼接写法返回 2 行（含 flag），参数化写法返回 0 行，正常标题仍可用（终端实录截图）}{web:param-fix}

\paragraph{结果与核对}
对照 \figref{web:param-fix} 的三行输出：\verb|fix BEFORE rows| 有两条记录、\verb|fix AFTER rows| 是空、\verb|normal title 'public' still works| 仍有记录。三行齐全才算修复验证通过。若你自己改了服务器做实验，记得把两组请求都留档。

\paragraph{为什么会成功}
机制一句话：\textbf{参数化查询把语句结构和数据分开提交}。数据库先按 \verb|?| 占位符把 SQL 编译成语法树，数据随后作为"一个值"填进去，永远不会参与语法分析。所以单引号、\verb|OR|、注释符在参数里都只是普通字符。反面教材是黑名单过滤（删掉 \verb|OR|、替换单引号）：大小写、双写、编码、注释变形都能绕过它。

\debugtip{① 注入串仍然生效：十有八九是服务没重启，跑的还是旧代码。② 正常标题也查不到了：占位符数量和参数个数不匹配，检查 \code{(title,)} 是不是写成了 \code{(title)}（后者不是元组）。③ 分不清测的是哪一版：把两版响应都存文件，文件名里写明 fixed / before。}

\paragraph{变式练习}
\begin{enumerate}
\item 用变式输入 \verb|public' OR '1'='1| 在修复版再测一次，确认同样只按字面值匹配。
\item 写一个"按标题搜索"的参数化小接口，再加一个按 \verb|body| 搜索的接口，想想两个参数是否都要参数化（答案：都要）。
\item 思考题：参数化能修 W1 的越权吗？为什么不能？（提示：W1 的病因是授权判断的位置，不是 SQL。）
\end{enumerate}

% ===========================================================================

\chapter{复盘：知识地图与易错清单}
\label{web:review}
\goal{把四道题的经验压成"下次遇到这类题怎么办"的固定流程。}

\section{Web 题五步流程}
\begin{enumerate}
\item \textbf{分诊}：这是 Web 题吗？给的是 URL、附件还是源码？
\item \textbf{建基线}：正常请求什么样子？状态码、响应头、正文各存一份。
\item \textbf{找可控输入}：URL 参数、POST 表单、Cookie、文件名、请求头，逐个列出。
\item \textbf{单变量试探}：一次改一个输入，把"数据被谁解释"猜出来（数据库 / 浏览器 / 文件系统 / 权限判断）。
\item \textbf{利用并复核}：构造输入拿到答案，再独立验证：对源码、重算、或平台反馈。
\end{enumerate}

\section{漏洞类型判断表}
\begin{center}
\begin{tabular}{p{0.2\textwidth}p{0.28\textwidth}p{0.22\textwidth}p{0.2\textwidth}}
\toprule
现象 & 可能的解释器 & 先做什么 & 修复方向 \\
\midrule
单引号/引号报 SQL 错 & 数据库 & 读拼接处源码 & 参数化查询 \\
条件恒真多返回数据 & 数据库 & 手工代入语句推结构 & 参数化查询 \\
输入出现在页面并被执行 & 浏览器 & 看输出位置 & 输出编码 \\
能读到目录外文件 & 文件系统 & 看路径怎么拼 & 规范化 + 根目录核对 \\
改个编号看到别人数据 & 权限判断 & 画权限小表 & 服务端核对归属 \\
改角色/隐藏域进后台 & 权限判断 & 试 POST/Cookie 是否也生效 & 会话角色判定 \\
\bottomrule
\end{tabular}
\end{center}

\section{易错清单}
\begin{enumerate}
\item 没有基线就改输入，最后说不清是哪个改动起的作用；
\item 把界面限制（隐藏按钮、置灰控件）当成权限检查；
\item 用 PowerShell 直敲 \verb|curl| 而不是 \verb|curl.exe|，被别名坑；
\item 引号嵌套写错，把 PowerShell 解析错误当成"SQL 报错"；
\item URL 手工编码出错（空格、引号、井号），改用 \verb|--data-urlencode|；
\item 只记 payload 不记机制，换个数据库方言就束手无策；
\item 修复后只测恶意输入，忘了确认正常功能还活着；
\item 用黑名单过滤当修复方案；
\item 改了服务端代码忘了重启，测的还是旧进程；
\item 看到 flag 就复制，不核对大小写、花括号和来源；
\item 对真实平台题硬编结论，把"题面说的功能"当成"已验证的漏洞"；
\item 把命令用在未经授权的系统上。
\end{enumerate}

\section{下次遇到这类题怎么办}
\textbf{给 URL 的题}：先抓一次正常请求（curl 或 F12），列出全部可控输入，按"数据库 / 浏览器 / 文件系统 / 权限"四选一逐个试。\textbf{给源码的题}：搜 \verb|SELECT|、\verb|open(|、\verb|eval|、\verb|execute|，追每个危险函数的参数来源。\textbf{给上传功能的题}：先走正常流程建基线，再动文件名、路径、类型。\textbf{权限题}：画"谁能看什么"的小表，逐格验证服务端是否核对。\textbf{卡住时}：汇报已知输入、观察到的响应、已排除的假设、下一条想验证的动作——而不是只说"做不出来"。

\memory{基线、单变量、机制解释、独立复核——四件事凑齐，才叫解题；只拿到 flag 叫猜中。}

% ===========================================================================
\chapter{自测与参考答案}
\label{web:answers}
\goal{先自测概念是否掌握，再核对完整推演。}

\section{自测题}
\begin{enumerate}
\item 一次 HTTP 请求由哪四段组成？空行的作用是什么？
\item GET 和 POST 的参数分别放在哪里？"放在 POST 里"为什么仍然不可信？
\item 把单引号写成 \verb|%27| 发送，数据库收到的是什么字符？
\item Cookie 由谁生成、由谁保存、由谁随请求发送？为什么它不是信任凭证？
\item SQL 的 \verb|SELECT ... WHERE title = 'public'| 里，哪部分是语句结构，哪部分是数据？
\item W2 的注入串里，\verb|--| 为什么要写？不写会发生什么？
\item SQL 注入、XSS、路径穿越分别由哪个组件把数据当成了代码？
\item 认证和授权的区别是什么？W1 的问题属于哪一个？
\item 为什么 W1 把参数从 GET 改成 POST 依然不安全？
\item 验证一次 SQL 注入修复，最少要做哪两组测试？
\end{enumerate}

\section{参考答案与完整推演}
\paragraph{第 1 题}请求 = 请求行（方法+路径+查询串+协议版本）+ 请求头 + 空行 + 请求体；响应 = 状态行 + 响应头 + 空行 + 响应体。空行把"头"和"体"分开，双方都遵守。

\paragraph{第 2 题}GET 参数在 URL 查询串里，POST 参数通常在请求体里。但请求整个由客户端构造，抓包或脚本都能自己写 POST 正文，所以"藏在 POST 里"只是不显眼，不是可信。

\paragraph{第 3 题}单引号。\verb|%27| 只是传输编码，服务器解码后还原，数据库看到的仍是 \verb|'|。

\paragraph{第 4 题}Cookie 由服务器通过响应头 \verb|Set-Cookie| 下发、浏览器保存、之后随每个请求自动带上。因为保存在客户端，其内容可被查看和篡改，所以携带 Cookie 不等于可信，服务端要用会话记录去验证它。

\paragraph{第 5 题}\verb|SELECT|、\verb|FROM|、\verb|WHERE title =| 和两端的引号是语句结构；\verb|public| 是数据。拼接的问题就是把外部数据写进了"结构"的位置。

\paragraph{第 6 题}\verb|--| 是 SQL 注释符，用来把原语句尾部剩下的引号和内容注释掉，保证语句语法完整。不写的话，输入拼出的语句尾部会多出孤立引号，多半以语法错误收场（可以自己在靶场验证）。

\paragraph{第 7 题}SQL 注入——数据库；XSS——浏览器；路径穿越——文件系统（再加上权限判断组件的越权问题）。

\paragraph{第 8 题}认证回答"你是谁"，授权回答"你能做什么"。W1 属于授权问题：服务端把客户端自报的角色字段当成了授权依据。

\paragraph{第 9 题}因为 POST 正文同样由客户端填写，把 \verb|role=admin| 放进请求体，服务器若照旧信任它，越权依旧。修复必须发生在服务端的授权判断上，与传输位置无关。

\paragraph{第 10 题}两组：① 恶意输入不再改变查询结构（注入失效）；② 正常输入照常工作（功能没被误伤）。\figref{web:param-fix} 就是两组一起做的例子。

\paragraph{完整推演：W1}
启动靶场后，\code{curl.exe -s -i} \url{http://127.0.0.1:8765/gate?role=guest} 返回 \verb|Access denied|。只把参数值改成 \verb|admin|，响应体变成 \code{flag\{http\_\allowbreak query\_\allowbreak is\_input\}}。对照源码（\figref{web:app-source}）：\verb|/gate| 分支直接比较 \verb|role| 与 \verb|"admin"|，没有任何会话或口令检查。修复：由服务端根据已认证的会话决定角色，客户端自报字段一律忽略。

\paragraph{完整推演：W2}
先用 \verb|public| 建基线（一条笔记）。发单引号得到 SQL 报错，说明引号进入了语句结构。源码显示拼接位置：
\begin{lstlisting}[language=Python]
sql = f"SELECT title, body FROM notes WHERE title = '{title}'"
\end{lstlisting}
手工代入 \verb|' OR 1=1 -- |，语句变成"标题为空，或 1=1，尾部引号被注释"，两行全返回，其中 staff 行是 \verb|flag{sql_parameters_matter}|。换成参数化写法后，同一输入只是一个完整标题值，返回 0 行；\verb|public| 仍正常。注意：这并不能修复 W1 的授权错误，两个入口要分别分析。

% ===========================================================================
\chapter{附录 A：平台题目复核与提交 checklist}
\label{web:platform}
\goal{把本地方法用于题目复核与组队交流，所有结论均以实际证据为准。}

\section{一眼区分附件题与在线环境题}
玄机首页 \url{https://xj.edisec.net/} 登录后进入"靶场"。平台题先把四件事记下：\textbf{题面、附件、环境、提交格式}，再选工具。提供可离线分析附件的题（如"五格电文"），第一步是下载附件、列文件、识别格式；提供在线环境的题（如题目页提供在线环境），第一步是读附件与题面，再从页面启动靶机并记录实际地址。题面中的 \verb|your_ip| 是占位写法。是否收费、按钮名称可能变化，以当前详情页为准。不要仅凭方向标签判断是否必须联网。

\section{组队复核要点}
\paragraph{T1：同时收到多份材料怎么分工？}先给每题写\textbf{交付物、第一条证据、需要的运行环境}，再指派负责人：URL 题先留正常请求和响应；密文先记录格式与长度；EXE 先静态识别打包方式与入口；ELF 先识别架构和保护；PCAP 先按会话列端点与协议。每题的关键证据都要有第二人复核。

\paragraph{T2：没有新证据时怎么汇报？}卡住时汇报四样：已知输入、观察到的响应、已经排除的假设、下一条想验证的动作。例如："已确认 guest 请求可以读取公开页面；下一步对照不同权限下同一请求的响应，并记录差异。"完整题解也应保留关键失败假设与适用条件。校内赛的目标是多人能复现同一证据链，再由平台确认最终答案。

\section{提交 checklist}
\begin{enumerate}[label=\arabic*.]
\item 题面、附件版本、环境地址已记录；附件原件保留，未覆盖。
\item 基线请求与改动后请求均已保存，能说清唯一变化的输入。
\item flag 有独立于最后打印语句的证据：W1 核对 \verb|labs/web/app.py| 对应分支；W2 确认记录来自本地服务数据库；平台题需再加密/重算或对照样本核对。
\item 提交前检查大小写、花括号和空格，按题面要求提交完整 flag。
\item 在平台页面等待提交完成，确认显示"FLAG 正确"及步骤完成后再记为已验证；未验证题记为"入口/阶段完成，答案未确认"，不编造。
\item 题解写清命令、脚本版本、附件路径、完整输出与平台反馈；不把账号口令写入讲义或共享包。
\end{enumerate}

% ===========================================================================
\chapter{附录 B：命令速查表}
\label{web:cheatsheet}
\goal{想做什么，就敲哪条命令。全部命令在 PowerShell 里执行，除非注明 WSL。}

\begin{longtable}{p{0.36\textwidth}p{0.56\textwidth}}
\toprule
想做什么 & 敲什么命令 \\
\midrule
\endhead
看 Python 版本 & \code{python --version} \\
启动本地靶场 & \code{python labs/web/app.py} \\
停止服务 & 在服务终端按 \code{Ctrl+C} \\
看完整响应（头+正文）& \code{curl.exe -s -i} \url{http://127.0.0.1:8765/} \\
请求门禁接口 & \code{curl.exe -s -i} \url{http://127.0.0.1:8765/gate?role=guest} \\
只改参数值做对照 & \code{curl.exe -s -i} \url{http://127.0.0.1:8765/gate?role=admin} \\
按标题搜索（自动编码）& \code{curl.exe -s -i -G --data-urlencode "title=public"} \url{http://127.0.0.1:8765/notes} \\
发单引号探测 & \code{curl.exe -s -G --data-urlencode 'title='''} \url{http://127.0.0.1:8765/notes} \\
发注入串 & \code{curl.exe -s -i -G --data-urlencode 'title='' OR 1=1 -- '} \url{http://127.0.0.1:8765/notes} \\
用 URL 编码直接访问 & \code{curl.exe -s -i} \url{http://127.0.0.1:8765/notes?title=%27%20OR%201%3D1%20--%20} \\
看响应头细节 & \code{curl.exe -v URL} \\
发 POST 表单 & \code{curl.exe -s -i -X POST -d "a=1" URL} \\
加自定义请求头 & \code{curl.exe -s -i -H "Cookie: k=v" URL} \\
确认 curl 是不是别名 & \code{Get-Alias curl | Format-List Name, ReferencedCommand} \\
看真正的 curl 版本 & \code{curl.exe --version} \\
看整个源码文件 & \code{Get-Content labs/web/app.py} \\
搜关键行带上下文 & \code{Select-String -Path labs/web/app.py -Pattern 'sql = f' -Context 1,2} \\
把响应存文件 & \code{curl.exe -s URL > out.txt} \\
看文件类型（WSL）& \code{file xxx} \\
开静态文件服务做对比 & \code{python -m http.server 8000} \\
跑 W3 对照演示脚本 & \code{python -X utf8} \url{figures/raw/fig-web-10-param-fix.py} \\
看当前目录（PowerShell/WSL）& \code{pwd} \\
列目录（PowerShell / WSL）& \code{dir} / \code{ls} \\
\bottomrule
\end{longtable}

% ===========================================================================
\chapter{附录 C：术语表}
\label{web:glossary}
\goal{按字母/拼音顺序速查，每个术语一句话解释。}

\begin{longtable}{p{0.26\textwidth}p{0.66\textwidth}}
\toprule
术语 & 一句话解释 \\
\midrule
\endhead
404 & 状态码，表示服务器上找不到请求的路径。 \\
Cookie & 服务器下发、浏览器保存并随请求自动携带的小段数据，常放会话标识。 \\
curl & 命令行 HTTP 客户端，把请求原样发出并打印响应；Windows 下要用 \code{curl.exe}。 \\
GET & HTTP 方法，取数据，参数通常放在 URL 查询串里。 \\
HTTP & 浏览器与网站之间"请求-响应"的通信协议。 \\
IP 地址 & 网络上一台机器的编号；\code{127.0.0.1} 特指"本机自己"。 \\
payload & 为触发漏洞而构造的输入数据，本讲义称"注入串"。 \\
POST & HTTP 方法，交数据，参数通常放在请求体里。 \\
SQL & 操作数据库的语言，用 \code{SELECT/FROM/WHERE} 等语句描述要什么数据。 \\
SQL 注入 & 用户输入被数据库当成 SQL 语法执行的漏洞。 \\
URL & 资源地址，形如 协议://主机:端口/路径?查询串。 \\
URL 编码 & 把特殊字符写成 \code{\%XX} 以便在 URL 中传输；解码后还原，不提供安全。 \\
XSS & 跨站脚本：用户输入被浏览器当成网页脚本执行的漏洞。 \\
靶场 & 专门练习用的、故意留漏洞的环境，通常只在本机或授权平台。 \\
参数 & URL 或表单中"名字=值"的一项，是客户端可控的输入。 \\
查询串 & URL 问号后面的部分，形如 \code{role=guest\&page=1}。 \\
单变量验证 & 一次只改一个输入的实验方法，用于确定因果。 \\
端口 & 同一台机器上区分不同服务的编号，本靶场用 8765。 \\
方法 & 请求行里的动作名，如 GET、POST。 \\
攻击面 & 系统中可被触碰、可被利用的入口集合。 \\
会话 & 服务器记录的"同一用户的一系列请求"，靠会话标识串联。 \\
基线 & 不做任何改动时的正常请求与响应，用于对比。 \\
路径 & URL 中主机后面的部分，指示要访问的功能或资源。 \\
路径穿越 & 借 \code{../} 一类写法访问允许目录之外文件的漏洞。 \\
请求 & 客户端发给服务器的一次"提问"。 \\
请求头 & 请求中"名字: 值"形式的元信息，如 Host、Cookie。 \\
请求体 & 请求中放数据的部分，POST 表单常放在这里。 \\
参数化查询 & SQL 写法：结构用占位符定死，数据单独提交，永不参与语法。 \\
权限 & 角色或身份被允许做的事情；由服务端判断才有意义。 \\
认证 & 核对身份（你是谁）的过程。 \\
状态码 & 响应开头的三位数字，表示这次请求的处理结果。 \\
注入 & 把数据写进"会被当成代码解释"的位置的攻击手法。 \\
越权 & 以别人的权限访问资源，或访问超出自己角色的功能。 \\
响应 & 服务器对一次请求的答复，含状态行、响应头与响应体。 \\
会话标识 & Cookie 中用来认出"同一登录用户"的随机串。 \\
占位符 & 参数化 SQL 里的 \code{?}，代表"这里将来放一个值"。 \\
\bottomrule
\end{longtable}

\end{document}
